\documentclass[11pt,a4paper]{letter}
\usepackage[margin=1in]{geometry}
\usepackage{hyperref}
\signature{Saeyong Park \\ Avikus Co., Ltd. \\ \href{mailto:saeyong.park@avikus.ai}{saeyong.park@avikus.ai}}
\address{Saeyong Park \\ Avikus Co., Ltd. \\ Seongnam, Republic of Korea}
\date{April 10, 2026}

\begin{document}
\begin{letter}{The Editor-in-Chief \\ \textit{Ocean Engineering} \\ Elsevier}

\opening{Dear Editor,}

Please find enclosed our manuscript titled \textbf{``WaveSplat:
Grid-Free Spectral Gaussian Wave Field Estimation from Sparse
Shipboard LiDAR''}, which we submit for consideration as a
research article in \textit{Ocean Engineering}.

The manuscript addresses a long-standing bottleneck in real-time
shipboard wave sensing: grid-based ensemble Kalman filters
(EnKF/ETKF), the de-facto standard for LiDAR and X-band wave field
reconstruction, must interpolate sparse, irregularly distributed
LiDAR returns onto a fixed spatial grid whose resolution is
independent of the true scene complexity. We propose
\textbf{WaveSplat}, a grid-free spectral representation in which the
wave field is a weighted sum of $N$ plane waves whose wavenumbers
and amplitudes are directly optimised against sparse observations
while the angular frequency is fixed by the linear dispersion
relation as a hard constraint. An energy-based pruning and a
residual-triggered growth mechanism make the number of components
\emph{adaptive}, matching scene complexity from one swell to
multi-component seas.

\textbf{Key contributions and findings:}
\begin{itemize}
  \item A grid-free spectral estimation framework with
        dispersion-constrained parameterisation and geometry-aware
        confidence weighting, runnable in real time on a single CPU
        thread.
  \item A pre-registered six-metric evaluation panel
        (nRMSE, variance ratio, Pearson $\rho$, Skill Score,
        directional $R^2$, $H_s$ error) on 100 episodes spanning
        Beaufort~2--5 and 1--30 spectral components, compared against
        a hyperparameter-tuned ETKF-grid (24-config sweep) and a
        sliding-window GPR.
  \item WaveSplat attains \textbf{30.1\% lower nRMSE}
        ($0.223$ vs.\ $0.319$) and runs $\sim\!10\times$ faster
        ($13.4$ vs.\ $132.6$~ms/frame) than the best ETKF
        configuration, and is the \emph{only} method with a positive
        Skill Score against a zero predictor.
  \item Improvements are highly significant over 100 paired runs
        (Wilcoxon $p{=}1.1\!\times\!10^{-15}$ vs.\ ETKF, paired
        Cohen's $d{=}1.16$; $p{=}4.7\!\times\!10^{-4}$ vs.\ sliding
        GPR, $d{=}0.42$), robust under two LiDAR raydrop models, and
        stable under 20--100\% observation-density subsampling.
\end{itemize}

\textbf{Relevance to Ocean Engineering.} The paper sits squarely in
the journal's scope of physical ocean measurement, data assimilation,
and model-based sensing for marine operations. It directly enables
downstream applications explicitly named in the journal's call,
including wave-preview feedforward control for high-speed hydrofoil
vessels (demonstrated in a companion paper), route optimisation, and
structural-health monitoring.

\textbf{Originality and prior communication.} This work is original,
has not been published previously, and is not under consideration
elsewhere. A short early version of the grid-free idea was
internally reviewed at Avikus but has not been submitted to any
journal or conference. No part of the manuscript has been previously
published.

\textbf{Data and reproducibility.} The source code of WaveSplat, the
physics-based LiDAR wave simulator (WaveSim), all experiment scripts,
and five representative episodes are openly available at
\url{https://github.com/SAEYONGPARK94/wavesplat}. The full 100-episode
dataset can be regenerated from the provided configuration files and
fixed random seeds, ensuring end-to-end reproducibility of every
table and figure.

\textbf{Suggested reviewers} (no conflict of interest with the
author):
\begin{itemize}
  \item Prof.\ Nicolas Desmars --- École Centrale Nantes ---
        X-band/LiDAR ETKF reconstruction.
  \item Prof.\ Marc Buckley --- Helmholtz-Zentrum Hereon ---
        Wave field reconstruction from remote sensing.
  \item Dr.\ Jonas Ehlers --- U-Net / FNO wave reconstruction.
  \item Prof.\ Ulrik Dam Nielsen --- DTU --- Ship-motion-based wave
        estimation.
\end{itemize}

The sole author declares no competing financial interests or
personal relationships that could have influenced this work. The
author has completed all listed CRediT roles
(conceptualisation, methodology, software, validation, formal
analysis, investigation, writing -- original draft, writing --
review \& editing, visualisation).

Thank you very much for considering our manuscript. We look forward
to the reviewers' comments.

\closing{Sincerely,}

\end{letter}
\end{document}
